\section*{Capítulo \ref{ch:g8:plastics} --- Los plásticos y los materiales sintéticos}

\begin{solution}{exo:g8:plastics:1}
Naturales: el algodón, la lana. Artificiales: el papel, la viscosa.
Sintéticos: el nailon, el poliestireno.
\end{solution}

\begin{solution}{exo:g8:plastics:2}
Del petróleo (o del gas natural).
\end{solution}

\begin{solution}{exo:g8:plastics:3}
El PET (poli(tereftalato de etileno)): botellas de agua y de refresco.
\end{solution}

\begin{solution}{exo:g8:plastics:4}
Un \omterm{def:g8:plastics:thermoplastic}{termoplástico} se ablanda cada vez que se calienta y se puede volver
a moldear; un \omterm{def:g8:plastics:thermoplastic}{termoestable} se endurece para siempre la primera vez y se
carboniza en lugar de ablandarse.
\end{solution}

\begin{solution}{exo:g8:plastics:5}
Polipropileno (PP), un \omterm{def:g8:plastics:thermoplastic}{termoplástico}: se puede fundir y volver a
moldear.
\end{solution}

\begin{solution}{exo:g8:plastics:6}
El triángulo solo indica el \omterm{def:g8:plastics:plastic}{plástico}, para ayudar a clasificarlo. Que el
objeto se recicle depende de que se recoja y de que haya una planta que
recicle ese \omterm{def:g8:plastics:plastic}{plástico}.
\end{solution}

\begin{solution}{exo:g8:plastics:7}
Una sartén se calienta: un \omterm{def:g8:plastics:thermoplastic}{termoplástico} se ablandaría, y un
\omterm{def:g8:plastics:thermoplastic}{termoestable} no.
\end{solution}

\begin{solution}{exo:g8:plastics:8}
$\frac{353}{460} \approx 0.77$: alrededor del \qty{77}{\percent}.
\end{solution}

\begin{solution}{exo:g8:plastics:9}
El PVC contiene cloro: al quemarlo desprende cloruro de hidrógeno, un
gas corrosivo, además del hollín y el monóxido de carbono de la
\omterm{def:g8:combustion-and-fuels:complete}{combustión incompleta} de una hoguera de jardín.
\end{solution}

\begin{solution}{exo:g8:plastics:10}
$300 \times 36 = 10\,800$ botellas; $10\,800 \times 25 = 270\,000$\,g, es
decir, \qty{270}{kg}.
\end{solution}

\begin{solution}{exo:g8:plastics:11}
Dióxido de carbono y agua. Su carbono viene del petróleo, almacenado bajo
tierra durante millones de años: al quemarlo se añade dióxido de carbono
al \omterm{def:g4:air-a-mixture-of-gases:air}{aire}, como al quemar un \omterm{def:g8:combustion-and-fuels:fossil}{combustible fósil}.
\end{solution}

\begin{solution}{exo:g8:plastics:12}
La botella no se pudre: el sol y las olas la rompen en trozos cada vez
más pequeños. Los trocitos flotan o se hunden en el agua, y los peces
se los tragan con la comida.
\end{solution}

\begin{solution}{pb:g8:plastics:1}
\textbf{1.} Los cuerpos: 1 (PET). Los tapones: 2 (PEAD) o 5 (PP).

\textbf{2.} Sintéticos: se fabrican a partir de sustancias construidas
por los químicos, casi siempre a partir del petróleo.

\textbf{3.} Sí: el PET es un \omterm{def:g8:plastics:thermoplastic}{termoplástico}.

\textbf{4.} Cada \omterm{def:g8:plastics:plastic}{plástico} se recicla por separado: fundidos juntos, una
\omterm{def:g6:pure-substances-and-mixtures:mixture}{mezcla} de \omterm{def:g8:plastics:plastic}{plásticos} da un material de mala calidad.

\textbf{5.} $0.80 \times 1000 = \qty{800}{kg}$.

\textbf{6.} Tapones: $0.12 \times 1000 = \qty{120}{kg}$. Etiquetas: $0.05
\times 1000 = \qty{50}{kg}$.

\textbf{7.} $0.03 \times 1000 = \qty{30}{kg}$ de suciedad y líquido.

\textbf{8.} $80 + 12 + 5 = \qty{97}{\percent}$.

\textbf{9.} $0.025 \times 800 = \qty{20}{kg}$.

\textbf{10.} Las escamas ahorran petróleo y la energía de fabricar PET
nuevo, y convierten un residuo en un objeto útil.

\textbf{11.} Dióxido de carbono (y, si la \omterm{def:g4:what-a-fire-needs:burning}{combustión} es mala, monóxido de
carbono y hollín).

\textbf{12.} $800 - 20 = \qty{780}{kg}$ de escamas de PET limpias por
bala.
\end{solution}
